% !TeX root = ../main.tex
\section{Khách hàng mục tiêu}

% % !TeX root = ../main.tex
% \subsection{Đối tượng khách hàng}

% \subsubsection{Người mới chơi eSports và sinh viên}

% Đây là nhóm người dùng có nhu cầu chơi game cao nhưng chưa có nhiều kinh nghiệm chọn gear hoặc có ngân sách hạn chế. Họ dễ bị ảnh hưởng bởi quảng cáo, review trên mạng hoặc lời khuyên từ cộng đồng, từ đó mua sản phẩm không thật sự phù hợp với nhu cầu cá nhân.

% Với nhóm này, Mutux giúp giảm rủi ro mua nhầm bằng cách cho phép thuê hoặc thử thiết bị trong một khoảng thời gian ngắn trước khi quyết định mua. Điều này đặc biệt phù hợp với sinh viên, vì họ thường không muốn bỏ ra 1--3 triệu đồng cho một món gear khi chưa chắc sản phẩm đó đáp ứng đúng nhu cầu.

% \subsubsection{Game thủ, đội tuyển, câu lạc bộ và phòng game}

% Nhóm này có nhu cầu tối ưu thiết bị để cải thiện trải nghiệm hoặc hiệu suất thi đấu. Họ quan tâm nhiều hơn đến độ phù hợp, độ ổn định và cảm giác sử dụng thực tế trong quá trình luyện tập hoặc thi đấu. Việc được thử nhiều mẫu gear khác nhau giúp họ đưa ra lựa chọn chính xác hơn trước khi mua.

% Đối với các câu lạc bộ eSports, đội tuyển hoặc phòng game, Mutux có thể hỗ trợ thuê thử nhiều thiết bị cho thành viên trước khi mua số lượng lớn. Phòng game cũng có thể dùng nền tảng để thử nghiệm các dòng gear mới trước khi trang bị cho hệ thống máy.

% \subsubsection{Shop gaming gear}

% Các cửa hàng bán thiết bị gaming có thể đưa sản phẩm lên nền tảng để người dùng thuê thử. Đây là một kênh marketing và bán hàng mới, giúp shop tăng độ tiếp cận, tạo thêm doanh thu từ thiết bị trưng bày hoặc thiết bị nhàn rỗi, đồng thời tăng khả năng chuyển đổi từ thuê thử sang mua thật.

\subsection{Đối tượng khách hàng}

Khách hàng mục tiêu của Mutux gồm hai nhóm chính trong mô hình marketplace hai chiều. Nhóm thứ nhất là \textbf{người thuê (Renter)}, bao gồm học sinh, sinh viên, game thủ, người đi làm, người sáng tạo nội dung và các đội tuyển hoặc câu lạc bộ eSports có nhu cầu trải nghiệm, sử dụng thiết bị gaming trong thời gian ngắn nhưng chưa muốn hoặc chưa đủ khả năng mua ngay. Nhóm thứ hai là \textbf{người cho thuê (Lender)}, bao gồm các cửa hàng gaming gear, phòng máy và cá nhân đang sở hữu thiết bị nhàn rỗi, mong muốn tạo thêm doanh thu từ tài sản hiện có.

Mutux tập trung kết nối hai nhóm khách hàng này thông qua dịch vụ thuê thử gaming gear, giúp người thuê giảm rủi ro mua nhầm và giúp người cho thuê tăng hiệu quả khai thác thiết bị. Sau khi xác định hai nhóm khách hàng tổng quát, dự án phân tích sâu hơn đặc điểm, nhu cầu và hành vi của từng nhóm thông qua hồ sơ khách hàng lý tưởng (Ideal Customer Profile - ICP) dưới đây.


% !TeX root = ../main.tex
\subsection{Chân dung khách hàng mục tiêu}

Để định hình chiến lược sản phẩm và các tính năng cốt lõi của \projectname{}, việc xác định rõ Chân dung khách hàng mục tiêu (Ideal Customer Profile - ICP) là vô cùng quan trọng. Do \projectname{} hoạt động theo mô hình nền tảng hai chiều (Double-sided Marketplace), ICP của dự án được phân tách rõ ràng thành hai nhóm đối tượng tương hỗ: \textbf{Người thuê (Renter)} và \textbf{Người cho thuê (Lender)}.

\subsubsection{Chân dung người thuê}

Nhóm người thuê trên \projectname{} được phân chia thành ba phân khúc khách hàng cốt lõi với các đặc điểm nhân khẩu học, tâm lý học và hành vi khác nhau, nhưng cùng chung mong muốn trải nghiệm thiết bị công nghệ cao cấp với chi phí tối ưu.

\paragraph{Phân khúc học sinh, sinh viên}
\begin{itemize}
    \item \textbf{Nhân khẩu học:} Độ tuổi từ 18 -- 23 tuổi. Học sinh THPT hoặc sinh viên các trường Đại học, Cao đẳng tại các đô thị lớn. Nguồn ngân sách cá nhân hoặc trợ cấp từ gia đình có giới hạn (trung bình từ 2 -- 4 triệu đồng/tháng).
    \item \textbf{Tâm lý học:} Có niềm đam mê lớn với game giải trí hoặc các bộ môn Esport (Valorant, League of Legends, FC Online, CS2, v.v.). Thích sở hữu và trải nghiệm các dòng gaming gear thời thượng để cải thiện hiệu suất chơi game hoặc thể hiện cá tính riêng, nhưng gặp rào cản lớn về tài chính.
    \item \textbf{Hành vi:} Thường xuyên theo dõi các kênh review công nghệ trên YouTube, TikTok và các cộng đồng mạng xã hội. Chơi game chủ yếu tại phòng trọ hoặc các phòng máy (cyber game). Thích mua sắm nhưng cân nhắc cực kỳ kỹ về giá, thường tìm kiếm các giải pháp tiết kiệm nhất.
    \item \textbf{Nỗi đau (Pain Points):}
    \begin{itemize}
        \item Ngân sách eo hẹp không đủ để chi trả từ 1,5 -- 3 triệu đồng cho một chiếc bàn phím cơ hay chuột gaming cao cấp.
        \item Nỗi sợ ``tiền mất tật mang'' khi mua nhầm thiết bị không phù hợp (ví dụ chuột quá to/nhỏ so với cỡ tay, phím cơ quá ồn ảnh hưởng đến bạn cùng phòng).
        \item Nhu cầu sử dụng tăng cao đột xuất nhưng ngắn hạn (ví dụ tham gia một giải đấu phong trào tại trường).
    \end{itemize}
    \item \textbf{Giải pháp:} Cung cấp dịch vụ thuê ngắn hạn theo ngày với chi phí chỉ bằng một phần nhỏ giá trị sản phẩm. Tích hợp tính năng hỗ trợ cọc linh hoạt để giảm tối đa gánh nặng tài chính ban đầu.
\end{itemize}

\paragraph{Phân khúc người đi làm, sáng tạo nội dung}
\begin{itemize}
    \item \textbf{Nhân khẩu học:} Độ tuổi từ 22 -- 28 tuổi. Lập trình viên, designer, video editor, hoặc freelancer tự do. Thu nhập ở mức trung bình khá (từ 8 -- 18 triệu đồng/tháng).
    \item \textbf{Tâm lý học:} Coi trọng chất lượng công cụ làm việc. Đòi hỏi sự thoải mái, hiệu năng và thẩm mỹ cao để tối ưu hóa năng suất lao động (ưa thích phím cơ custom, chuột công thái học, tai nghe cách âm tốt). Quan tâm đến xu hướng setup bàn làm việc cá nhân (desk setup).
    \item \textbf{Hành vi:} Làm việc với máy tính liên tục từ 8 -- 10 tiếng mỗi ngày tại nhà hoặc văn phòng. Thích tự mình trải nghiệm trước khi đưa ra quyết định nâng cấp thiết bị. Thường tham khảo các đánh giá từ cộng đồng chuyên sâu.
    \item \textbf{Nỗi đau (Pain Points):}
    \begin{itemize}
        \item Rủi ro bị chấn thương cổ tay (RSI) hoặc mệt mỏi khi chọn nhầm chuột/bàn phím không đạt tiêu chuẩn công thái học.
        \item Việc đầu tư một bộ thiết bị cao cấp (tổng chi phí từ 5 -- 10 triệu đồng) cần sự chắc chắn về trải nghiệm thực tế lâu dài, điều mà các cửa hàng bán lẻ không thể đáp ứng chỉ qua vài phút dùng thử tại showroom.
        \item Không có nhiều thời gian rảnh rỗi để di chuyển đến nhiều cửa hàng khác nhau để trải nghiệm thử.
    \end{itemize}
    \item \textbf{Giải pháp:} Cho phép thuê trải nghiệm sâu từ 3 -- 7 ngày ngay tại không gian làm việc thực tế của họ với dịch vụ giao nhận tận nơi tiện lợi, giúp người dùng tự tin đưa ra quyết định mua sắm cuối cùng.
\end{itemize}

\paragraph{Phân khúc game thủ bán chuyên}
\begin{itemize}
    \item \textbf{Nhân khẩu học:} Độ tuổi từ 18 -- 25 tuổi. Thành viên các đội tuyển Esport bán chuyên hoặc các câu lạc bộ Esport của các trường Đại học.
    \item \textbf{Tâm lý học:} Có tư duy thi đấu nghiêm túc, đặt nặng tính thắng thua và thành tích. Đòi hỏi cực kỳ khắt khe về các thông số kỹ thuật của thiết bị (mắt đọc cảm biến chuột nhạy, tần số phản hồi cao, switch bàn phím tốc độ cực nhanh, tai nghe hỗ trợ định vị âm thanh chuẩn xác).
    \item \textbf{Hành vi:} Luyện tập cường độ cao (4 -- 6 tiếng/ngày) trước các giải đấu. Thường xuyên di chuyển thi đấu giữa các phòng máy hoặc địa điểm tổ chức sự kiện.
    \item \textbf{Nỗi đau (Pain Points):}
    \begin{itemize}
        \item Thiết bị cá nhân bị hỏng hóc hoặc gặp sự cố bất ngờ ngay trước thềm giải đấu quan trọng.
        \item Không đủ kinh phí để đầu tư các thiết bị chuẩn thi đấu Esport quốc tế (thường có giá rất đắt đỏ) chỉ để phục vụ cho một giải đấu ngắn ngày kéo dài từ 3 -- 5 ngày.
    \end{itemize}
    \item \textbf{Giải pháp:} Cung cấp các gói thuê thiết bị thi đấu cao cấp (Logitech G Pro, Razer Ultimate, v.v.) trong thời gian diễn ra giải đấu với tốc độ xử lý đơn hàng và giao nhận nhanh chóng, đảm bảo thiết bị đạt trạng thái tốt nhất.
\end{itemize}

\paragraph{Bảng tổng hợp người thuê}

\begin{table}[H]
    \centering
    \caption{So sánh các phân khúc Người thuê (Renter) trên nền tảng}
    \label{tab:renter-icp}
    \small
    \begin{tabular}{p{3.5cm}p{3.5cm}p{3.5cm}p{3.5cm}}
        \toprule
        \textbf{Tiêu chí} & \textbf{Game thủ học sinh, sinh viên} & \textbf{Lập trình viên, Creator trẻ} & \textbf{Game thủ bán chuyên, thi đấu} \\
        \midrule
        \textbf{Độ tuổi} & 18 -- 23 tuổi & 22 -- 28 tuổi & 18 -- 25 tuổi \\
        \textbf{Khả năng chi trả} & Thấp (Dưới 3tr/tháng) & Trung bình khá (8 -- 18tr/tháng) & Thấp đến trung bình \\
        \textbf{Mục đích thuê} & Trải nghiệm gear xịn giải trí ngắn ngày & Thử nghiệm công thái học phục vụ làm việc & Thiết bị chuyên dụng Esport cho giải đấu \\
        \textbf{Nỗi đau chính} & Thiếu kinh phí mua mới; Sợ mua nhầm gear không hợp & Rủi ro do chọn sai thiết bị làm việc lâu dài, không có nhiều thời gian & Thiết bị cá nhân hỏng bất ngờ; Thiếu gear chuẩn thi đấu \\
        \textbf{Giải pháp: } & Giá thuê rẻ; Hỗ trợ đặt cọc linh hoạt & Thuê trải nghiệm thực tế 3 -- 7 ngày tại nhà & Gói thuê ngắn ngày chuẩn Esport; Giao nhận cực nhanh \\
        \bottomrule
    \end{tabular}
\end{table}

\subsubsection{Chân dung người cho thuê}

Để đảm bảo nguồn cung đa dạng và chất lượng trên sàn \projectname{}, dự án hướng tới ba phân khúc đối tác cho thuê chính.

\paragraph{Phân khúc cửa hàng kinh doanh vừa và nhỏ}
\begin{itemize}
    \item \textbf{Đặc điểm doanh nghiệp:} Các shop kinh doanh gaming gear độc lập, đại lý phân phối quy mô vừa và nhỏ tại địa phương. Sở hữu một lượng hàng trưng bày (demo units) và hàng cũ đã qua sử dụng nhất định trong cửa hàng.
    \item \textbf{Tâm lý học:} Mong muốn gia tăng doanh thu bán hàng, tối ưu hóa các tài sản sẵn có và cạnh tranh hiệu quả với các chuỗi bán lẻ khổng lồ bằng dịch vụ chăm sóc khách hàng và trải nghiệm thực tế độc đáo.
    \item \textbf{Hành vi:} Quảng bá sản phẩm thông qua các hội nhóm mạng xã hội và phục vụ khách vãng lai tại địa bàn. Thường xuyên gặp tình trạng hàng trưng bày bị lỗi thời và mất giá trị theo thời gian.
    \item \textbf{Nỗi đau (Pain Points):}
    \begin{itemize}
        \item Hàng trưng bày (demo) bám bụi tại showroom và khấu hao liên tục nhưng không trực tiếp tạo ra dòng tiền.
        \item Khó tiếp cận tệp khách hàng trực tuyến mới do hạn chế về ngân sách marketing và độ nhận diện thương hiệu thấp hơn các chuỗi lớn.
        \item Tỷ lệ chuyển đổi khách hàng từ xem sang mua còn thấp do khách hàng ngần ngại trước các sản phẩm giá trị cao.
    \end{itemize}
    \item \textbf{Giải pháp:} Tạo ra một kênh cho thuê chuyên nghiệp, giúp shop tối ưu hóa hiệu suất khai thác hàng trưng bày/hàng cũ. Đồng thời, đóng vai trò như một phễu marketing hiệu quả, chuyển đổi người thuê thành người mua sản phẩm mới trong tương lai.
\end{itemize}

\paragraph{Phân khúc cá nhân sở hữu thiết bị nhàn rỗi}
\begin{itemize}
    \item \textbf{Nhân khẩu học:} Độ tuổi từ 20 -- 32 tuổi. Có thu nhập ổn định và đam mê sâu sắc với bộ môn gaming gear hoặc bàn phím cơ custom.
    \item \textbf{Tâm lý học:} Thích sưu tầm, thay đổi và trải nghiệm liên tục các dòng thiết bị mới. Tự hào về bộ sưu tập cá nhân và muốn chia sẻ niềm đam mê công nghệ với cộng đồng. Mong muốn tối ưu hóa nguồn vốn nhàn rỗi để tiếp tục nuôi dưỡng đam mê nâng cấp gear.
    \item \textbf{Hành vi:} Sở hữu đồng thời nhiều thiết bị (3 -- 5 chiếc bàn phím cơ custom, nhiều dòng chuột gaming cao cấp). Hoạt động rất tích cực trên các diễn đàn công nghệ hoặc hội nhóm bàn phím cơ.
    \item \textbf{Nỗi đau (Pain Points):}
    \begin{itemize}
        \item Các bộ gear đắt tiền để không tại nhà lâu ngày dễ bị ẩm mốc, hỏng hóc hoặc lỗi thời mà bán lại thì chịu lỗ lớn.
        \item Việc tự cho thuê cá nhân qua các kênh mạng xã hội tự phát tiềm ẩn rủi ro rất cao (bị lừa đảo, quỵt cọc, tráo đổi linh kiện chính hãng bên trong thiết bị).
        \item Thiếu quy trình giao dịch, hợp đồng pháp lý và quy chuẩn bàn giao rõ ràng để bảo vệ tài sản của bản thân.
    \end{itemize}
    \item \textbf{Giải pháp:} Cung cấp hệ thống giao dịch trung gian an toàn tuyệt đối với quy trình xác thực KYC nhằm xác minh danh tính người dùng nghiêm ngặt, thanh toán ký quỹ (escrow) tạm giữ tiền cọc, và bộ quy chuẩn minh chứng bàn giao (proof) rõ ràng dưới sự hỗ trợ giải quyết tranh chấp từ Admin.
\end{itemize}

\paragraph{Phân khúc chủ quán net/gaming center}
\begin{itemize}
    \item \textbf{Đặc điểm:} Các chủ quán net hoặc gaming center quy mô nhỏ và vừa, đang sở hữu nhiều máy tính, chuột, bàn phím, tai nghe, tay cầm hoặc thiết bị dự phòng phục vụ khách hàng.
    \item \textbf{Nhu cầu và động lực:} Muốn khai thác thêm doanh thu từ thiết bị ít sử dụng, thu hút game thủ và hỗ trợ các đội tuyển, câu lạc bộ eSports hoặc giải đấu tổ chức tại quán.
    \item \textbf{Nỗi đau (Pain Points):} Thiết bị dự phòng bị khấu hao nhưng chưa tạo thêm giá trị; chủ quán không muốn tự quản lý việc đăng tin, đặt cọc, giao nhận và xử lý hư hỏng ngoài hoạt động kinh doanh chính.
    \item \textbf{Giải pháp:} Mutux cung cấp quy trình cho thuê, đặt cọc và kiểm tra thiết bị tập trung. Quán có thể bắt đầu với một số phụ kiện phù hợp, áp dụng hình thức nhận--trả tại quán để giảm chi phí vận hành và rủi ro giao nhận.
\end{itemize}


\begin{table}[H]
    \centering
    \caption{So sánh các phân khúc Người cho thuê (Lender) trên nền tảng}
    \label{tab:lender-icp}
    \small
    \begin{tabular}{p{3.2cm}p{3.2cm}p{3.2cm}p{3.2cm}}
        \toprule
        \textbf{Tiêu chí} & \textbf{Cửa hàng gaming gear} & \textbf{Chủ quán net/gaming center} & \textbf{Cá nhân sưu tầm gear} \\
        \midrule
        \textbf{Đặc điểm nguồn cung} & Hàng trưng bày (demo units), sản phẩm cũ hoặc đã qua sử dụng, số lượng lớn và đa dạng. & Máy và phụ kiện đang phục vụ khách, kèm thiết bị dự phòng hoặc ít sử dụng. & Bộ sưu tập cá nhân, số lượng ít nhưng giá trị cao và độc đáo. \\
        \textbf{Động cơ chính} & Tối ưu hóa hàng mẫu; tạo phễu tiếp thị; chuyển đổi người thuê thành người mua. & Tạo thêm doanh thu từ thiết bị nhàn rỗi và thu hút khách hàng đến quán. & Tạo thu nhập thụ động để duy trì đam mê nâng cấp gear. \\
        \textbf{Nỗi đau lớn nhất} & Hàng trưng bày bị khấu hao mà không sinh lợi; khó tiếp cận khách hàng online mới. & Không muốn tự quản lý đặt cọc, giao nhận, vệ sinh và rủi ro hư hỏng. & Sợ bùng cọc, mất cắp hoặc tráo linh kiện khi giao dịch tự phát. \\
        \textbf{Giải pháp \projectname{}} & Cung cấp kênh tiếp thị và cho thuê thử trước khi mua. & Hỗ trợ đăng cho thuê, đặt cọc, kiểm tra và nhận--trả tại quán. & KYC, escrow, đối soát minh chứng (proof) và giải quyết tranh chấp. \\
        \bottomrule
    \end{tabular}
\end{table}

\subsubsection{Phân khúc trọng tâm giai đoạn MVP}

Để đảm bảo hiệu quả vận hành và tối ưu hóa nguồn lực có hạn trong giai đoạn đầu, \projectname{} xác định chiến lược tiếp cận trọng tâm như sau:
\begin{itemize}
    \item \textbf{Về phía Người thuê (Renter):} Nhóm \textbf{Game thủ học sinh, sinh viên (Casual \& Student Gamers)} tại các khu vực đô thị và làng đại học lớn (ví dụ: Khu đô thị ĐHQG-HCM) sẽ là tệp khách hàng mục tiêu cốt lõi. Đây là nhóm có nhu cầu trải nghiệm gear rất cao, mật độ tập trung đông và khả năng lan tỏa thông tin tự nhiên (word-of-mouth) tốt nhất, giúp nền tảng nhanh chóng đạt được lượng người dùng hoạt động tối thiểu (critical mass).
    \item \textbf{Về phía Người cho thuê (Lender):} Các \textbf{Cửa hàng Gaming Gear vừa và nhỏ (SME Gear Shops)} địa phương sẽ là đối tác cung cấp nguồn cung ưu tiên. Hợp tác với các shop giúp \projectname{} giải quyết bài toán đa dạng hóa danh mục sản phẩm nhanh chóng, đảm bảo chất lượng kỹ thuật của thiết bị bàn giao và tận dụng các chính sách vận chuyển sẵn có của shop, giảm thiểu rủi ro vận hành so với việc nhận gear từ cá nhân riêng lẻ ở giai đoạn đầu.
\end{itemize}
Sự kết hợp này tạo ra một vòng tuần hoàn cung -- cầu khép kín, an toàn và nhanh chóng, làm tiền đề vững chắc trước khi mở rộng dịch vụ sang nhóm lập trình viên, người sáng tạo nội dung và các cá nhân có gear nhàn rỗi trong cộng đồng.


% !TeX root = ../main.tex
\subsection{Chân dung Khách hàng điển hình}

\subsubsection{Chân dung Người thuê thiết bị}

Hai persona này tập trung vào nhóm \textbf{Game thủ học sinh, sinh viên (Casual \& Student Gamers)} tại các khu vực đô thị và làng đại học lớn, đặc biệt là những cụm có mật độ sinh viên cao như Khu đô thị ĐHQG-HCM. Đây là tệp người thuê cốt lõi của Mutux trong giai đoạn đầu vì có nhu cầu trải nghiệm gaming gear cao, ngân sách mua mới còn hạn chế, thường sinh hoạt trong các cộng đồng trường/lớp/ký túc xá và có khả năng lan tỏa tự nhiên qua bạn bè, câu lạc bộ, phòng máy hoặc nhóm Discord. Việc tập trung vào nhóm này giúp nền tảng dễ đạt lượng người dùng hoạt động tối thiểu (critical mass) hơn so với việc tiếp cận nhiều phân khúc rời rạc ngay từ đầu.

\paragraph{Persona 1: Minh Quân -- Sinh viên muốn trải nghiệm gaming gear tốt, mới ra mắt}

\begin{itemize}
    \item \textbf{Nhân khẩu học:} Nam, 20 tuổi, sinh viên năm hai ngành Công nghệ thông tin tại Trường Đại học Công nghệ Thông tin, ĐHQG-HCM. Sống tại ký túc xá ĐHQG, ngân sách cá nhân khoảng 2--4 triệu đồng/tháng, thường chi tiêu cho ăn uống, đi lại, game và một phần nhỏ cho thiết bị học tập/giải trí.
    \item \textbf{Bối cảnh sử dụng:} Quân chơi Valorant, CS2 và các game co-op cùng nhóm bạn 4--5 buổi vào tối hàng tuần tại ký túc xá. Quân muốn dùng thử chuột nhẹ, bàn phím cơ, tai nghe tốt hơn thiết bị phổ thông đang có và các mẫu gaming gear mới ra mắt để chơi thoải mái hơn, bắt kịp xu hướng trong nhóm bạn và biết sản phẩm nào hợp với mình trước khi mua.
    \item \textbf{Nỗi đau:} Không đủ ngân sách mua ngay một bộ gear 3--6 triệu đồng; các sản phẩm mới ra mắt thường có giá cao, vượt quá khả năng chi trả của sinh viên; sợ mua nhầm chuột không hợp tay hoặc switch bàn phím quá ồn trong phòng trọ; chỉ muốn trải nghiệm trong vài ngày/cuối tuần nhưng thị trường chủ yếu bán mới; lo ngại thuê/mượn qua người lạ vì không rõ tình trạng thiết bị và trách nhiệm khi hỏng hóc.
    \item \textbf{Động lực:} Muốn nâng trải nghiệm chơi game và setup cá nhân mà không phải đầu tư lớn; muốn thử thiết bị thực tế vài ngày trước khi quyết định mua; thích được trải nghiệm sớm sản phẩm mới và đổi mới gear theo xu hướng trong cộng đồng bạn bè. Vì ngân sách hạn chế, Quân đặc biệt quan tâm đến mô hình cọc linh hoạt, chẳng hạn cọc trả sau hoặc giảm áp lực cọc ban đầu cho người thuê có lịch sử tốt.
    \item \textbf{Hành vi tiếp cận:} Dễ gặp nhất khi chơi game chung, qua vòng tròn bạn bè, nhóm Discord nhỏ, Facebook group trường, ký túc xá và lời giới thiệu từ người quen đã từng thuê. Đây là persona phù hợp nhất cho MVP vì có mật độ tập trung cao, phản hồi nhanh và khả năng tạo word-of-mouth trong cộng đồng sinh viên.
\end{itemize}

\paragraph{Persona 2: Nhật Khang -- Game thủ có thi đấu giải}

\begin{itemize}
    \item \textbf{Nhân khẩu học:} Nam, 22 tuổi, sinh viên năm cuối ngành Công nghệ thông tin tại Trường Đại học Khoa học Tự nhiên, ĐHQG-HCM, thường tham gia các giải Valorant/CS2 sinh viên và giải cộng đồng bán chuyên. Có đội chơi cố định, tự góp ngân sách cá nhân hoặc quỹ đội để đăng ký giải, thuê phòng máy và chuẩn bị thiết bị.
    \item \textbf{Bối cảnh sử dụng:} Nhu cầu của Khang tăng mạnh theo mùa giải. Trước các giải kéo dài 2--5 ngày, Khang cần chuột nhẹ, bàn phím phản hồi nhanh và tai nghe định vị âm thanh tốt để luyện tập và thi đấu ổn định. Ngoài mùa giải, nhu cầu thuê thấp hơn; trong mùa giải, thời gian chuẩn bị ngắn khiến việc có gear đúng mẫu, đúng thời điểm trở nên quan trọng.
    \item \textbf{Nỗi đau:} Mua gear chuẩn esports quá tốn kém nếu chỉ dùng nhiều trong một mùa giải hoặc một giải ngắn ngày; khó tiếp cận thiết bị đúng chuẩn thi đấu với chi phí hợp lý; mượn thiết bị qua người quen không ổn định và thiếu cam kết về thời gian giao nhận; cần giao nhận nhanh, đúng mẫu thiết bị và tình trạng kỹ thuật rõ ràng.
    \item \textbf{Động lực:} Muốn thuê ngắn ngày hoặc theo mùa giải thiết bị chuẩn esports, độ trễ thấp và tình trạng tốt; sẵn sàng trả phí cao hơn thuê thông thường nếu thiết bị giúp chuẩn bị thi đấu ổn định; quan tâm đến quy trình đặt cọc, bàn giao và hoàn trả minh bạch; có động lực giới thiệu dịch vụ cho đồng đội hoặc đội khác nếu Mutux đáp ứng tốt giai đoạn cao điểm mùa giải.
    \item \textbf{Hành vi tiếp cận:} Dễ gặp qua lịch giải đấu sinh viên, phòng máy, câu lạc bộ esports, admin cộng đồng game và đối tác tổ chức sự kiện. Persona này có mức độ nỗi đau rất cao theo mùa giải, phù hợp để \projectname{} thử nghiệm gói thuê ngắn ngày/theo mùa cho game thủ sinh viên có thi đấu giải và tạo trường hợp tham chiếu trong cộng đồng esports.
\end{itemize}

Nhìn chung, hai persona người thuê đều thuộc cùng một cụm khách hàng mục tiêu ban đầu: học sinh, sinh viên chơi game tại các khu vực tập trung đông người trẻ. Họ khác nhau về cường độ sử dụng và bối cảnh thuê, nhưng cùng chia sẻ nhu cầu trải nghiệm gear tốt với chi phí thấp hơn mua mới, khả năng phản hồi nhanh và khả năng giới thiệu dịch vụ cho bạn bè. Vì vậy, thay vì mở rộng ngay sang nhiều nhóm người thuê rời rạc, Mutux nên ưu tiên đạt độ phủ trong các cộng đồng trường đại học, ký túc xá, phòng máy và câu lạc bộ esports để tạo critical mass cho MVP.

% \begin{table}[H]
%     \centering
%     \caption{So sánh hai chân dung người thuê điển hình}
%     \label{tab:renter-personas-pa3}
%     \small
%     \begin{tabular}{p{0.22\textwidth}p{0.34\textwidth}p{0.34\textwidth}}
%         \toprule
%         \textbf{Tiêu chí} & \textbf{Minh Quân -- Sinh viên} & \textbf{Nhật Khang -- Game thủ thi đấu giải} \\
%         \midrule
%         Bối cảnh sử dụng & Thuê gear ngắn ngày để chơi cùng nhóm, trải nghiệm setup tốt hơn và thử trước khi mua. & Thuê gear chuẩn thi đấu trước và trong giải. \\
%         Mức độ nỗi đau & Cao do ngân sách mua mới hạn chế và sợ chọn sai thiết bị. & Rất cao trước giải; thiết bị lỗi có thể ảnh hưởng trực tiếp đến thành tích đội. \\
%         Khả năng chi trả & Trả được đơn nhỏ theo ngày; có thể chịu ảnh hưởng từ bạn bè/gia đình khi mua mới. & Tự trả hoặc góp cùng đội/CLB; ưu tiên gói thuê ngắn ngày. \\
%         Kênh tiếp cận & CLB esports, nhóm trường, ký túc xá, Discord, phòng máy. & Giải đấu sinh viên, phòng máy, CLB esports, admin cộng đồng game. \\
%         Vai trò với MVP & Tạo giao dịch đầu, phản hồi nhanh và lan truyền trong cộng đồng sinh viên. & Kiểm chứng nhu cầu thuê mùa giải và tạo trường hợp tham chiếu trong cộng đồng esports. \\
%         \bottomrule
%     \end{tabular}
% \end{table}

\subsubsection{Chân dung Người cho thuê Thiết bị}

Phía cung của Mutux không nên được hiểu là mọi cá nhân đang sở hữu một bộ gear đều sẵn sàng cho thuê. Khi các thiết bị phổ biến ngày càng dễ mua và giá thuê của một món đồ không còn cách quá xa giá mua, động lực cho thuê của cá nhân sẽ thấp; trong khi đó, chi phí kiểm tra, vệ sinh, giao nhận và rủi ro hư hỏng vẫn tồn tại. Vì vậy, ở giai đoạn đầu, Mutux ưu tiên các đối tác đã có sẵn thiết bị, không gian và tệp khách hàng. Với họ, cho thuê là một dịch vụ bổ sung để tăng trải nghiệm hoặc doanh thu, thay vì trở thành nguồn thu chính.

\paragraph{Persona 3: Nam -- Chủ quán net/gaming center}

\begin{itemize}
    \item \textbf{Bối cảnh:} Nam vận hành một quán net hoặc gaming center quy mô nhỏ đến vừa quanh khu vực trường đại học, khu dân cư và các điểm tập trung game thủ. Quán đã có sẵn chuột, bàn phím, tai nghe, tay cầm hoặc một vài bộ máy dự phòng để phục vụ khách và các nhóm chơi theo giờ.
    \item \textbf{Nhu cầu và vấn đề:} Doanh thu chính vẫn đến từ giờ chơi, nhưng thiết bị dự phòng hoặc phụ kiện ít được sử dụng có thể bị khấu hao mà chưa tạo thêm giá trị. Nam cũng muốn giữ khách bằng trải nghiệm khác biệt, song không muốn tự xây dựng một quy trình cho thuê riêng, phải theo dõi đặt cọc, tình trạng thiết bị và giao nhận ngoài hoạt động của quán.
    \item \textbf{Động lực tham gia:} Có thêm doanh thu từ những thiết bị sẵn có; cho phép khách thuê nhanh phụ kiện khi quên đồ hoặc muốn thử một setup tốt hơn; hỗ trợ các đội game, câu lạc bộ và giải đấu nhỏ tổ chức tại quán. Mutux đóng vai trò mang thêm khách và xử lý phần đặt thuê, còn quán giữ quyền quyết định thiết bị nào được đưa lên hệ thống.
    \item \textbf{Rào cản:} Lo ngại thất lạc, hư hỏng, vệ sinh thiết bị và việc cho thuê làm ảnh hưởng đến số máy phục vụ khách tại giờ cao điểm. Do đó, Nam phù hợp với mô hình thử nghiệm phạm vi hẹp: vài phụ kiện có nhu cầu rõ ràng, nhận--trả tại quán và quy định bồi thường đơn giản.
    \item \textbf{Hành vi tiếp cận:} Có thể tiếp cận qua các quán net gần HCMUS, Thủ Đức, câu lạc bộ esports, admin cộng đồng game và đơn vị tổ chức giải sinh viên. Đây là persona cung cấp phù hợp vì đã có sẵn tài sản, nhân sự hỗ trợ và địa điểm nhận--trả.
\end{itemize}

\paragraph{Persona 4: Bình -- Chủ shop gear quy mô nhỏ}

\begin{itemize}
    \item \textbf{Bối cảnh:} Bình điều hành một shop gear độc lập có showroom nhỏ hoặc kết hợp bán hàng online. Shop bán chuột, bàn phím, tai nghe, màn hình và phụ kiện gaming; một phần hàng được dùng để trưng bày, cho khách trải nghiệm, hoặc là hàng cũ/đổi trả còn sử dụng tốt. Các showroom gear vốn đã có vai trò cho khách xem và trải nghiệm sản phẩm trực tiếp, như mô hình showroom của GEARVN tại Thành phố Thủ Đức \cite{gearvnShowroom}.
    \item \textbf{Nhu cầu và vấn đề:} Biên lợi nhuận của các món phổ thông có thể không cao, còn khách hàng thường muốn thử cảm giác cầm chuột, độ nảy bàn phím hoặc chất âm trước khi mua. Tuy nhiên, Bình không muốn mọi sản phẩm trong kho đều phải chuyển sang cho thuê vì sẽ phát sinh kiểm kê, vệ sinh, hao mòn và rủi ro làm giảm giá trị bán lại.
    \item \textbf{Động lực tham gia:} Dùng Mutux như một kênh ``dùng thử trước khi mua'', biến hàng demo hoặc một số sản phẩm chọn lọc thành doanh thu phụ, đồng thời đưa khách thuê quay lại shop để mua phụ kiện phù hợp. Mục tiêu chính của Bình vẫn là bán hàng; khoản phí thuê và dữ liệu phản hồi chỉ là phần bổ trợ giúp tăng khả năng chốt đơn.
    \item \textbf{Rào cản:} Cần quy định rõ thiết bị nào được cho thuê, mức đặt cọc, trách nhiệm khi trầy xước và cách đối soát doanh thu. Shop chỉ nên đưa lên Mutux các thiết bị có giá trị trải nghiệm cao hoặc khó quyết định khi mua, thay vì các món giá rẻ mà khách có thể mua ngay.
    \item \textbf{Hành vi tiếp cận:} Phù hợp để tiếp cận trực tiếp qua các shop địa phương, cộng đồng gear, sự kiện gaming và đối tác tổ chức giải. Một gói thử nghiệm đơn giản (một vài sản phẩm, thời hạn ngắn, nhận--trả tại shop) sẽ giúp Bình đánh giá nhu cầu mà không phải thay đổi mô hình bán hàng hiện tại.
\end{itemize}


Hai persona trên cho thấy nguồn cung khả thi của Mutux nên bắt đầu từ đối tác kinh doanh có sẵn hoạt động vận hành, thay vì kỳ vọng vào người dùng cá nhân. Quán net tạo ra điểm nhận--trả và nhu cầu thuê tại chỗ; shop gear tạo ra điểm trải nghiệm và cơ hội chuyển đổi sang mua hàng. Các cá nhân sở hữu gear có thể được mở rộng ở giai đoạn sau, khi Mutux đã chứng minh quy trình đặt cọc, kiểm tra và xử lý hư hỏng đủ đơn giản để họ cảm thấy an toàn.


